\documentclass[preamble=development-preamble.tex]{markdown}
title: A direct functional interpretation of arithmetic R in PTS-rk
date: October 20, 2026

:::Abstract
This development interprets the original functional presentation of
Tarski--Mostowski--Robinson arithmetic R in the unchanged, annotated
rank-based pure type system PTS-rk. Arithmetic expressions have a fixed
Church representation with zero, successor, addition and multiplication.
Certified observations supply equality and order, finite observations
separate numerals, and one positive bounded-case invariant establishes
the quantified axiom schemes. The formalization proves preservation of
an independently defined classical first-order deduction system, complete
proof/refutation results for ground atoms, and failure of faithfulness for
arbitrary sentences. Constructive enumeration of finite typing derivations
also gives executable interfaces returning actual annotated proof terms
with proofs of their original PTS typing judgments.
:::

# The result and its scope

The library is `Interpretation_of_R.functional`. Its public entry point is
`DirectRTheorem7.v`, and its principal declaration is `Theorem7.theorem7`.
This is the result labelled `thm:direct-R` in the harmonized arithmetic
manuscript, whose printed numbering gives Theorem 7. The construction is
independent of the relational interpretation and of the set-model libraries.
It shares structural lemmas, negative logic, finite-selector infrastructure
and certificate search with the parent `Interpretation_of_R` library.

Write $A^\dagger$ for `FirstOrder.sentence_translation A`, and
$\vdash p:A$ for the original `Typing.typ` judgment with
`Ranks.RankSorts`, natural-number sorts and the empty context. Every proof term below is
an element of `PTSrk.Syntax.Term nat 0`; open constructions use its
scope-indexed counterpart. The abbreviation `Realizes A` means that such
a term exists. It is not a new typing system or an assumed realization
relation.

:::Theorem
The declaration `Theorem7.theorem7` proves the conjunction of four claims.

1. Every sentence recognized by `FunctionalSyntax.axiom` has a closed
   originally typed inhabitant of its translation.
2. If `FunctionalSyntax.derives [] A`, then
   $\exists p.\;\vdash p:A^\dagger$.
3. For every ground equality or order atom, its numerical truth gives
   an inhabitant of its translation, and its numerical falsity gives
   an inhabitant of the translated atom's negation.
4. Some sentence has an inhabited translation but is not derivable in
   the independent source deduction system.
:::

The last clause is intentional: this is a deduction-preserving
interpretation with exact ground proof/refutation behavior, not a
conservativity theorem for all first-order sentences. Internal refutation
means an actual term of type $A^\dagger\to\bot$. External uninhabitation
of $\bot$ is a separate metatheoretic question and is not a conjunct of
`theorem7`.

# The unchanged target calculus and logical encoding

PTSrk has sorts $U_i$ indexed by external natural numbers, with
$U_i:U_0$ and product rule

$$
\frac{\Gamma\vdash A:U_a\qquad\Gamma,x:A\vdash B:U_b}
 {\Gamma\vdash\Pi x:A.B:U_{\max(a+1,b)}}.
$$

Abstractions retain their domain annotations. Conversion is the original
beta conversion; neither cumulativity nor eta conversion is used. Ranks
in this document are exact sort indices of displayed types, not implicit
universe inclusions.

The proof encodings use the fixed answer type

$$
\bot:=\Pi Z:U_0.Z:U_1,
\qquad \top:=U_0\to U_0:U_1.
$$

Set $\neg_c A=A\to\bot$ and
$A\wedge_c B=(A\to B\to\bot)\to\bot$. A flat finite disjunction is

$$
\bigvee_c(A_0,\ldots,A_k)
 :=(A_0\to\bot)\to\cdots\to(A_k\to\bot)\to\bot.
$$

A proof type $A$ is stable when $(\neg_c\neg_c A)\to A$ is inhabited.
Negative atoms and the translated connectives are stable, and dependent
products preserve stability of their codomain. This supports the classical
rules needed by the source calculus while restricting elimination of
negative encodings to stable targets. Positive domain evidence requires
the separate construction described below.

Exact-rank lifts handle motives whose natural rank is too small. With
$J_1=\top$, $J_{k+1}=J_k\to J_k$, and canonical inhabitants $j_k$,
the type $L_k(A)=J_{k-1}\to A$ has rank $k$ whenever $A:U_a$,
$a\leq k$ and $k\geq2$. Abstraction and evaluation at $j_{k-1}$
give maps between $A$ and its lift; no cumulativity rule is introduced.

# Arithmetic presentations and the admitted domain

The carrier in `Codes.Domain` is

$$
D:=\Pi X:U_1.\;X\to(X\to X)\to(X\to X\to X)
                  \to(X\to X\to X)\to X:U_3.
$$

For $\xi=(X,z,s,a,m)$ write
$\operatorname{ev}_{\xi}(d)=d\,X\,z\,s\,a\,m$. The notation
$\lambda\xi$ abbreviates five annotated abstractions, not a record
constructor. The four closed operations are

$$
\begin{aligned}
0_D&=\lambda\xi.z,&
S_Dd&=\lambda\xi.s(\operatorname{ev}_{\xi}(d)),\\
d+_De&=\lambda\xi.a(\operatorname{ev}_{\xi}(d))
                         (\operatorname{ev}_{\xi}(e)),&
d\mathbin{\cdot_D}e&=\lambda\xi.m(\operatorname{ev}_{\xi}(d))
                         (\operatorname{ev}_{\xi}(e)).
\end{aligned}
$$

Thus addition and multiplication construct arithmetic expression trees.
Their ground arithmetic equations will hold in interpreted equality;
they are not beta computations to canonical numerals. Put
$\nu_n=S_D^n0_D$ for the canonical presentation of external $n$.

For $Q:D\to U_{10}$, let $\mathsf{Cl}_S(Q)$ express successor closure
and $\mathsf{Cl}_f(Q)$ express binary closure under $f$. Define

$$
\mathsf{Good}(d):=
\Pi Q:D\to U_{10}.\;Q(0_D)\to\mathsf{Cl}_S(Q)
 \to\mathsf{Cl}_{+_D}(Q)\to\mathsf{Cl}_{\cdot_D}(Q)\to Q(d).
$$

This rank-12 predicate is the fixed admitted domain. Its constructor
certificates apply the corresponding closure argument to the children's
certificates. There is no claim of induction for arbitrary raw inhabitants
of $D$, nor is the desired bounded-case theorem inserted into the domain
definition. The rank-10 motive is fixed independently of every axiom index.

`DirectRCodes.v` constructs these expressions using the shared named
assembler, then closes them into original de Bruijn terms. Assembly proves
that names are resolved, not that terms are well typed.
`DirectRFoundations.v` separately proves their original typing, often by
the sound native checker. `FunctionalTerms.term_typed` and
`FunctionalTerms.term_admitted` extend typing and admission compositionally
to all source arithmetic terms under a typed, admitted environment.

# The certified observation telescope

An observation has eight data parameters:

$$
X:U_1,\quad z:X,\quad s:X\to X,\quad a,m:X\to X\to X,
\quad\delta:X\to U_4,\quad R_0,C_0:X\to X\to U_1.
$$

Its local proof relations are $u\simeq v=R_0(u,v)\to\bot$ and
$u\preceq v=C_0(u,v)\to\bot$. The data are followed by 23 separate
proof parameters. Every object variable in a law is followed immediately
by its $\delta$ guard. The certificates comprise:

- four domain laws: zero and closure under the three operations;
- equivalence of $\simeq$, congruence for successor, addition and
  multiplication, and congruence of $\preceq$ in both arguments;
- reflexivity and transitivity of $\preceq$;
- five arithmetic equations and five order propagation laws below.

With $1_X=sz$, the arithmetic and propagation laws are

$$
\begin{aligned}
a(z,u)&\simeq u,& a(su,v)&\simeq s(a(u,v)),\\
m(z,u)&\simeq z,& m(su,v)&\simeq a(v,m(u,v)),\\
m(u,z)&\simeq z,
\end{aligned}
$$

$$
\begin{gathered}
u\preceq su,\qquad u\preceq a(u,v),\qquad v\preceq a(u,v),\\
1_X\preceq v\;\longrightarrow\;u\preceq m(u,v),\qquad
1_X\preceq u\;\longrightarrow\;v\preceq m(u,v).
\end{gathered}
$$

All certificates have rank at most five. In `Codes.observation_fields`
they are literal successive binders, so a universal observation proof
must accept every one of them. They are parameters of a definable proof
type, not axioms added to PTS-rk.

Writing $\Pi\mathcal A$ for this complete telescope, define

$$
\begin{aligned}
E(d,e)&:=\Pi\mathcal A.\;
       \operatorname{ev}_{\mathcal A}(d)\simeq
       \operatorname{ev}_{\mathcal A}(e),\\
O(d,e)&:=\Pi\mathcal A.\;
       \operatorname{ev}_{\mathcal A}(d)\preceq
       \operatorname{ev}_{\mathcal A}(e),\\
G(d)&:=\Pi\mathcal A.\;\delta(\operatorname{ev}_{\mathcal A}(d)).
\end{aligned}
$$

These are `Codes.Equality`, `Codes.Order` and `Codes.G`, all with
rank-six values. $E$ and $O$ are stable; $G$ contains positive domain
evidence. Fold equations prove constructor closure of $G$. Instantiating
$\mathsf{Good}$ with the rank-10 lift of $G$ gives
$\mathsf{Good}(d)\to G(d)$. Opening the telescope and supplying these
guards transfers its equations, congruences and preorder laws to
observed arithmetic presentations.

External induction on $m,n$ now produces originally typed ground
addition and multiplication proofs, exported by `Kernel.R1_realized`
and `Kernel.R2_realized`. This induction constructs target terms;
native arithmetic automation only settles their numerical side conditions.

# Finite observations and internal separation

For an external cutoff $K$, put $c=K+1$ and use the chain
$\{0,\ldots,c\}$ with capped operations

$$
s(i)=\min(i+1,c),\quad
a(i,j)=\min(i+j,c),\quad m(i,j)=\min(ij,c).
$$

`DirectRArithmetic.v` establishes the numerical table properties.
Their target implementation uses $q=K+2$ rank-one selectors

$$
F_q:=\underbrace{U_1\to\cdots\to U_1}_{q\text{ arguments}}\to U_1,
\qquad e_i=\lambda Z_0\cdots Z_{q-1}.Z_i.
$$

A unary table routes a selector to its chosen output; a binary table
uses a selector for its row followed by a selector for its column.
For relation tables, a true entry returns $\bot$ and a false entry
returns $\top$. Negating the result gives an inhabited relation in
the true case and a refutable relation in the false case.

The selector-domain certificate is

$$
\delta_q(u):=\Pi Q:F_q\to U_3.\;
Q(e_0)\to\cdots\to Q(e_{q-1})\to Q(u):U_4.
$$

Its elimination rule reduces a guarded law to the finitely many selector
cases. This certifies every observation field, including binary closure,
transitivity, congruence and the multiplication laws.
`FiniteSpecialization.contextual_source_observation` assembles the
full observation; `FunctionalEvaluation.finite_numeral_evaluation`
relates numeral folds to their selectors. No assertion that every
raw inhabitant of $F_q$ is a selector is needed.

:::Theorem
`FunctionalSeparation.no_back_realized` constructs, for every $n$, a
closed proof of
$$O(\nu_{n+1},\nu_n)\to\bot.$$
More generally, reversed numeral comparisons and distinct-numeral
equalities have internal refutations.
:::

:::Proof
Choose cutoff $K=n$. The relevant folds are the distinct selectors
$e_{n+1}$ and $e_n$. Apply a proposed global order proof to the entire
closed finite observation. Its conclusion reduces to $\top\to\bot$;
application to the identity inhabitant of $\top$ yields $\bot$.
Ground order, transitivity and equality congruence extend this argument
to arbitrary reversed pairs and unequal numerals.
:::

The finite observations are separating target constructions. A single
finite chain is not asserted to be a model of all source numeral-separation
axioms, and numerical table validation alone is not the target proof.

# The positive bounded-case invariant

For a fixed external $n$, set

$$
\mathsf{Case}_n(d):=
\bigvee_c\bigl(E(d,\nu_0),\ldots,E(d,\nu_n),O(\nu_{n+1},d)\bigr).
$$

The last alternative is a global lower bound on $d$, rather than only
overflow in one observation. This distinction makes the case lemma strong
enough for both quantified arithmetic schemes.

:::Theorem
`FunctionalCases.bounded_case_realized` proves, for every $n$,
$$\Pi d:D.\;\mathsf{Good}(d)\to\mathsf{Case}_n(d).$$
:::

:::Proof
Use the rank-10 positive continuation package
$$
\mathsf{Inv}_n(d):=\Pi Z:U_8.\;(G(d)\to\mathsf{Case}_n(d)\to Z)\to Z.
$$
Prove closure under all four constructors. Open the parent's continuation,
then eliminate each child invariant at that same $Z$. This supplies
positive $G$ evidence and the negative case proof separately. Construct
the parent's $G$ evidence using the domain closure laws. To construct its
case proof, first open the case continuations, so all remaining branches
return $\bot$.

Zero takes the reflexive equality branch. Successor sends $i<n$ to
$i+1$; the $i=n$ branch becomes the large branch. A large child remains
large by $O(d,S_Dd)$. Addition combines two finite cases by their
ground sum; if that sum exceeds $n$, ground comparison supplies the large
branch. A large summand bounds the sum below by order propagation.

Multiplication needs a separate zero analysis. Two finite cases give
their ground product. A zero factor gives the zero branch using the
corresponding left or right zero law. A nonzero finite factor is at
least one; a large factor is also at least one because $1\leq n+1$.
Hence a large factor paired with a nonzero factor bounds the product
below by the appropriate guarded multiplication law. These arguments
cover finite/large, large/finite and large/large pairs.

Apply the parent continuation to the resulting $G$ and case proofs.
The lemma `good_invariant_realized` in `FunctionalCases` instantiates $\mathsf{Good}$
at this proved constructor-closed motive. Finally choose
$Z=\mathsf{Case}_n(d)$ and continuation $\lambda g\,h.h$.
:::

`DirectRFunctionalBounded.v` proves zero and successor closure;
`DirectRFunctionalAddition.v` and `DirectRFunctionalMultiplication.v`
prove the binary steps. Keeping both multiplication zero laws and the
positive-factor premises is essential: an unconditional claim that a
factor bounds its product is false when the other factor is zero.

The uniform rank accounting is:

| Construction | Rank |
| --- | --- |
| Arithmetic carrier $D$ | 3 |
| Finite selector domain $\delta_q(u)$ | 4 |
| Individual observation certificates | at most 5 |
| $E(d,e)$, $O(d,e)$ and $G(d)$ | 6 |
| $\mathsf{Case}_n(d)$ | 8 |
| $\mathsf{Inv}_n(d)$ | 10 |
| $\mathsf{Good}(d)$ | 12 |

These ranks do not grow with the cutoff or axiom index. No minimality
of the chosen ranks is claimed.

# Realizing all five source schemes

The functional source language has $0,S,+,\cdot,=,\leq$. For external
$m,n$, its schemes are the ground addition and multiplication equations,
distinctness of unequal numerals, and

$$
\begin{aligned}
\text{R4: }&\forall x.\;x\leq\bar n\to
                    \bigvee_{i=0}^{n}x=\bar i,\\
\text{R5: }&\forall x.\;x\leq\bar n\lor\bar n\leq x.
\end{aligned}
$$

`FunctionalSyntax.axiom` independently specifies these schemes. In
particular, functional R4 is the displayed implication. The source
definition does not import the relational presentation's biconditional.

`FunctionalAxioms.R4_realized` combines bounded cases with an assumed
$O(d,\nu_n)$. Each finite branch is already a permitted equality; the
large branch composes to $O(\nu_{n+1},\nu_n)$ and is internally refuted.
For `FunctionalAxioms.R5_realized`, finite branches give $O(d,\nu_n)$
by ground comparison and congruence. The large branch gives
$O(\nu_n,d)$ by transitivity through $\nu_{n+1}$.

Together with ground equations and separation, these give
`all_axioms_realized` in `FunctionalAxioms` for the concrete axiom-type
interface. The lemma `source_axiom_realized` in `FunctionalSourceAxioms` connects that
result to the independently defined source syntax and its translation.
Both are uniform in their indices, rather than a finite suite of examples.

# First-order translation and preservation of deduction

`FirstOrder.term` and `FirstOrder.formula` are indexed by their number
of free variables. Formulas have primitive equality, order, falsity,
implication, conjunction, disjunction and both quantifiers. Terms
translate compositionally to the four operations. Atoms use $E,O$;
implication uses an arrow and the other connectives use the negative
encodings. Quantifiers are relativized to admission:

$$
\begin{aligned}
(\forall x.A)^\dagger
 &=\Pi x:D.\;\mathsf{Good}(x)\to A^\dagger,\\
(\exists x.A)^\dagger
 &=\bigl(\Pi x:D.\;\mathsf{Good}(x)\to A^\dagger\to\bot\bigr)\to\bot.
\end{aligned}
$$

The independent `FunctionalSyntax.derives` includes hypotheses, theory
axioms, introduction and elimination for the primitive connectives,
classical double-negation elimination, equality reflexivity and equality
replacement. Its quantifier rules enforce the eigenvariable conditions
through scope indices and explicit weakening of the assumption context.
No constructor defines source derivability by target inhabitation.

Three layers make the induction on source derivations possible.

1. In `FunctionalTranslation`, `instantiate_translation` and the
   surrounding renaming/substitution lemmas connect source instantiation
   with actual target substitution.
2. In `FunctionalFormulaLogic`, `formula_formed` and `formula_stable`
   establish the formation and classical elimination obligations for
   every formula.
3. In `FunctionalEquality`, `formula_transport` proves equality transport
   between two typed, admitted environments that are pointwise $E$-related.
   Its quantified case extends both environments with the same admitted
   variable, preserving guards under binders.

This last condition matters: equality replacement is proved on admitted
arguments. It does not require $\mathsf{Good}$ to be invariant under
$E$ on all raw $D$-objects. The proof does not obtain admission of one
raw term merely from equality to another.

:::Theorem
`FunctionalDeduction.preservation` translates an arbitrary derivation
under assumptions into target inhabitation under any typed, admitted
environment realizing those assumptions. Its closed specialization,
`FunctionalDeduction.sentence_preservation`, states
$$
\mathsf{derives}([],A)\quad\Longrightarrow\quad
\exists p.\;\vdash p:A^\dagger.
$$
:::

The proof follows the source derivation. Stable elimination handles
conjunction projections, disjunction and existential elimination, and
the classical rule. Quantifier introduction extends the target context
with both an object and its admission proof; instantiation uses term
admission and the substitution theorem. Equality replacement invokes
the general formula-transport result. All final judgments are the
original `Typing.typ` judgments.

# Ground completeness and a precise nonfaithfulness witness

For a closed arithmetic term $t$, `FunctionalTerms.value` computes its
external natural value. `FunctionalTerms.ground_evaluation` constructs
$E(t^\dagger,\nu_{\mathsf{value}(t)})$ by structural induction, using
constructor congruence and the ground equations. Ground truth then
reduces to numeral equality or order; false cases use the finite
separation proofs. `Theorem7.ground_correctness` packages this for
`Theorem7.equality_atom` and `Theorem7.order_atom`.

For nonfaithfulness, take $\rho=\forall x.\;x\leq x$.
`Theorem7.order_reflexivity_realized` proves $\rho^\dagger$: admission
gives $G(d)$, which supplies the guard for each observation's reflexivity
certificate.

`FunctionalCountermodel` independently proves that $\rho$ is
underivable. Its carrier is $\mathbb N\cup\{\omega\}$, represented
by `option nat`. Successor fixes $\omega$; addition and multiplication
have their usual values on natural inputs and return zero whenever
an input is $\omega$. Order is ordinary order on naturals, every natural
is below $\omega$, and no ordered pair with left endpoint $\omega$
holds. Thus all five source schemes hold, whereas
$\omega\not\leq\omega$.

The formal semantics uses stable negative interpretations for logical
connectives. Decidability of atomic equality and order gives atomic
stability; structural stability then validates the classical source rule
constructively. `FunctionalCountermodel.soundness` proves soundness of
every source derivation for this semantics, and
`FunctionalCountermodel.order_reflexivity_underivable` obtains the
contradiction at $\omega$. Together these establish
`Theorem7.nonfaithfulness` without a native classical axiom. This is a
source countermodel, not a set-theoretic model of PTS-rk.

# Effective proof terms and their computational limitations

An existential proof in Rocq's `Prop` is not itself an executable
projection of its witness into `Type`. The development closes that gap
using two shared modules, `DirectRDerivationCertificates.v` and
`DirectRWitnessSearch.v`.

Finite tagged trees encode applications of exactly the original typing
rules. Conversion certificates give finite forward beta paths to a
common reduct; the original confluence characterization makes these
certificates complete. `DerivationCertificates.derive_sound` proves
that a successfully decoded certificate yields an original typing
derivation, and `DerivationCertificates.derive_complete` proves that
every original derivation has such a certificate.

`WitnessSearch.trees` enumerates finite sets with increasing bounds on
tree depth and labels. Every certificate occurs at some stage. Searching
a stage is decidable and returns a term only if its derived type is
syntactically the requested type. Existence of an inhabitant, together
with certificate completeness, proves that some stage succeeds.
Constructive unbounded search over natural numbers therefore produces

$$
\mathsf{witness}(A):\mathsf{Realizes}(A)\to
     \{p:\mathsf{Term}\;\mathsf{nat}\;0\mid\vdash p:A\}.
$$

The Type-valued public interfaces in `Theorem7` are `effective_axiom`,
`effective_source_axiom`, `effective_deduction`, `effective_ground_true`
and `effective_ground_false`. Each returns a term paired with its
original typing proof. In particular,
`Theorem7.theorem5_target : Kernel.Theorem5Target` supplies a function
selecting a typed witness for every concrete axiom instance. The
historical record name does not identify a different source theory.

No choice axiom, typing-decidability assumption or fixed sufficient fuel
bound is used. The procedure is intentionally exhaustive and has no
practical complexity guarantee. For small explicit certificates the
native checker is the appropriate computational interface. The older
optional generator in `DirectRGenerated.v` and its named totality/soundness
contracts in `DirectRObligations.v` are separate: completion of those
contracts is not a premise of the effective interfaces above.

# Source organization and reproduction

The source groups follow the mathematical dependencies:

- `DirectRCodes.v`, `DirectRFoundations.v` and `DirectRProofs.v` define
  the fixed annotated objects, certify their typing and prove ground laws.
- The `DirectRFinite` modules, together with the functional finite-law,
  evaluation and separation modules, implement and specialize the whole
  finite observation telescope.
- The functional bounded, addition, multiplication, cases and axioms
  modules prove the uniform invariant and all-index realization.
- The functional syntax, translation, formula-logic, source-axiom,
  equality and deduction modules connect independent source proofs
  to target inhabitants.
- `DirectRFunctionalCountermodel.v` and `DirectRTheorem7.v` supply
  nonfaithfulness and the final existential and effective interfaces.
  `DirectRFunctionalAudit.v` prints their types and transitive assumptions.

Use the repository's pinned environment: OCaml 4.14.2, Dune 3.23.1,
Rocq core 9.2.0, Stdlib 9.1.0 and Equations 1.3.2+9.2. The setup is
documented in `docs/DirectR-CI.md`. From the repository root:

```sh
dune build @theories/Interpretation_of_R/functional/audit
```

The following import exposes the public interfaces without importing
the relational or set-model libraries:

```coq
From Interpretation_of_R.functional
  Require Import DirectRTheorem7.
Check Theorem7.theorem7.
Check Theorem7.theorem5_target.
Check Theorem7.effective_deduction.
Print Assumptions Theorem7.theorem7.
Print Assumptions Theorem7.theorem5_target.
Print Assumptions Theorem7.effective_deduction.
```

Use a fresh absolute Dune build directory when a new assumption report
is required; cached success does not rerun printing commands. The
full local check is `bash tools/check-local-ci.sh`, including actionlint,
the staged developments and an isolated pristine upstream build.

The retained theories-layout verification of 14 October 2026 records a
fresh complete build and public reports showing `theorem7`,
`theorem5_target` and `effective_deduction` closed under the global
context. Its commands and raw results are in the
`theories-layout-2026-10-14` subdirectory of the handoff's verification
archive. Consult `docs/direct-r-handoff/STATE.md` for the latest
repository-wide verification. This documentation does not turn a cached build into a new
kernel audit. Historical draft headers and manuscript status sentences
describe their original checkpoints; current public declarations and
live verification determine the theorem's present status.

# Discussion

The functional construction keeps a fixed interface while allowing
unboundedly large finite proofs. The expression carrier exposes all four
constructors; the observation telescope supplies just enough algebra
to reason about them. Finite selectors provide genuine target proofs of
separation, and the bounded invariant transports the positive evidence
needed to use those algebraic laws through negative case analysis.

This also explains the interpretation's boundaries. It does not claim
that every raw carrier term is admitted, that observational equality
coincides with beta conversion, or that arbitrary source theorems can be
recovered from target proofs. Uniformly bounded ranks concern the
encoding, not the size or running time of the externally constructed
witnesses. The checked result establishes the full stated interpretation
and its constructive witness interface; rank optimization and an efficient
proof compiler are separate possible improvements.
